\section{Sonstiges}
Hier finden Sie einige letzte Hinweise zur Erstellung Ihrer Abschlussarbeit.

\subsection{Silbentrennung}
Achten Sie darauf, dass die automatische Silbentrennung in Ihrem Dokument aktiviert ist. Sollte sie nicht wirksam sein, prüfen Sie die Einstellungen Ihrer Textverarbeitung bzw. LaTeX-Konfiguration.

\subsection{Abspeichern der Arbeit}
Speichern Sie Ihre Arbeit regelmäßig und sichern Sie Zwischenstände. So können Sie im Bedarfsfall auf ältere Versionen zurückgreifen. Bei einem umfangreichen Anhang empfiehlt sich ein separates Anhangsverzeichnis zur besseren Orientierung.

\subsection{Platzprobleme}
Falls Sie Schwierigkeiten haben, die vorgegebene Seitenbegrenzung einzuhalten, prüfen Sie das Layout Ihrer Arbeit kritisch. Überlegen Sie, ob sich Inhalte in den Anhang auslagern lassen. Achten Sie außerdem darauf, keine unnötigen Seitenumbrüche in Überschriften zu verwenden.

\subsection{Rechtschreibprüfung}
Nutzen Sie eine Rechtschreibprüfung, um vermeidbare Tippfehler zu eliminieren. Verlassen Sie sich jedoch nicht ausschließlich auf automatische Korrekturen, da diese nicht immer zuverlässig sind.

\subsection{Einsatz von generativen IT-/KI-Werkzeugen}

Die Nutzung generativer KI-Tools wie ChatGPT, Bard, Dall-E oder Stable Diffusion kann Studierende bei der Ideenfindung, beim Strukturieren von Inhalten oder zur sprachlichen Überarbeitung unterstützen. Ebenso können solche Werkzeuge zur Generierung von Bildmaterial oder zur Hilfe bei der Texterstellung im technischen Kontext verwendet werden. Dabei dient die KI ausschließlich als Hilfsmittel zur Unterstützung des wissenschaftlichen Arbeitsprozesses.

\bigskip
\noindent
\textbf{Wichtig:} Gemäß den Richtlinien der Hochschule müssen alle verwendeten KI-Werkzeuge vollständig in einer eigenen Rubrik \textit{„Übersicht verwendeter Hilfsmittel“} aufgeführt werden. Dies beinhaltet:
\begin{itemize}
  \item den Produktnamen (z.\,B. \textit{ChatGPT}),
  \item die Zugriffsquelle (z.\,B. URL),
  \item eine Beschreibung der genutzten Funktionen,
  \item sowie den Umfang der Nutzung.
\end{itemize}

Wörtliche oder sinngemäße Übernahmen von KI-generierten Inhalten sind analog zu anderen Quellen zu kennzeichnen.
